\chapter{Kontrasepsi}\label{ch:g8:contraception}

Bab yang lalu memetakan sebuah mekanisme yang berkekuatan
menakjubkan: dua \omterm{def:g5:first-look-at-cells:cell}{sel} bertemu, dan empat puluh pekan kemudian
ada seorang manusia. Manusia adalah satu-satunya \omterm{def:g5:biodiversity:species}{spesies}
yang memahami mekanismenya sendiri --- dan pemahaman membawa
serta sebuah pilihan yang tak pernah dihadapi kodoknya:
apakah dan kapan menjalankannya. \emph{\omterm{prop:g8:contraception:principle}{Kontrasepsi}} adalah
ilmu terapan pilihan itu, dan ia bekerja, setiap caranya,
dengan menyela sebuah langkah yang kini dapat kamu sebutkan.

\section{Asasnya: sebutkan langkahnya, hadanglah langkahnya}

\begin{proposition}[Kontrasepsi dipetakan pada mekanismenya]\label{prop:g8:contraception:principle}
\omterm{def:g5:human-reproduction:pregnancy}{Kehamilan} membutuhkan sebuah rantai yang tak terputus
(\cref{prop:g8:sexual-reproduction:core},
\cref{prop:g8:from-egg-to-birth:firstweek}): \omterm{def:g8:sexual-reproduction:gametes}{gametnya}
dihasilkan --- spermanya disimpankan --- renang ke \omterm{def:g8:reproductive-systems:female}{saluran
telurnya} --- sebuah \omterm{def:g4:animal-reproduction:sexual}{sel telur} dilepaskan lalu ditemui ---
\omterm{def:g5:flower-to-fruit:steps}{pembuahan} --- penanaman. Yang disebut
\emph{kontrasepsi}\index{kontrasepsi} adalah setiap
penyelaan yang disengaja atas rantai ini. Setiap cara di rak
apotekernya adalah serangan atas satu mata rantai yang
bernama, dan memilahnya menurut mata rantainya itulah
seluruh teorinya:
\begin{itemize}
  \item cara \emph{penghadang} menghentikan pertemuan \omterm{def:g8:sexual-reproduction:gametes}{gametnya};
  \item cara \emph{hormonal} menghentikan \omterm{def:g4:animal-reproduction:sexual}{sel telurnya}
        dilepaskan sama sekali;
  \item cara lain membuat pertemuan atau penanamannya
        gagal.
\end{itemize}
\end{proposition}
\admitted

\section{Cara utamanya, menurut mata rantainya}

\begin{definition}[Cara penghadang: kondom]\label{def:g8:contraception:condom}
Sebuah \emph{kondom}\index{kondom} --- sarung tipis yang
dipakai pada \omterm{def:g8:reproductive-systems:male}{penis} (versi perempuannya melapisi vagina) ---
menghadang rantainya pada langkah penyimpanannya: air
maninya tidak pernah masuk ke dalam tubuh pasangannya.
Tersedia luas, dibutuhkan hanya ketika dibutuhkan --- dan
istimewa di raknya karena sebuah jasa kedua: hanya cara
inilah yang sekaligus menghadang lewatnya \emph{infeksi}
yang ditularkan lewat hubungan seksual, pokok dunia
\cref{ch:g9:microbes-and-infection}. Tugas ganda itulah
sebabnya kondom tetap ada dalam gambaran bahkan ketika cara
lain sedang dipakai.
\end{definition}

\begin{definition}[Cara hormonal: pilnya]\label{def:g8:contraception:pill}
Sebuah \emph{pil kontrasepsi}\index{pil kontrasepsi} memakai
permesinan \cref{def:g8:puberty:hormones} sendiri: diminum
setiap hari, \omterm{def:g8:puberty:hormones}{hormonnya} menahan daur \omterm{def:g8:reproductive-systems:female}{indung telurnya} dalam
keadaan tertunda (peristiwa 4
\cref{prop:g8:reproductive-systems:cycle} tanpa \omterm{def:g5:human-reproduction:pregnancy}{kehamilan}
apa pun) --- \emph{tanpa \omterm{prop:g8:reproductive-systems:cycle}{ovulasi}, tanpa \omterm{def:g4:animal-reproduction:sexual}{sel telur}, tak ada
yang dapat dibuahi}. Diresepkan oleh dokter, yang
mencocokkan caranya dengan orangnya; sangat andal bila
diminum teratur --- dan tidak menawarkan apa pun terhadap
infeksi, karena itulah kemitraan tetap \omterm{def:g8:contraception:condom}{kondomnya}.
\end{definition}

\begin{example}[Raknya, dipilah menurut mata rantainya]\label{ex:g8:contraception:shelf}
Cara umum lainnya, diarsipkan menurut mata rantai yang
diserangnya: susuk dan koyo hormonal --- asas pilnya dengan
ingatan yang lebih baik (berbulan-bulan atau
bertahun-tahun, tanpa minum harian); alat kecil yang
dipasang dokter di dalam \omterm{prop:g5:human-reproduction:plan}{rahim} --- pertemuan dan
penanamannya dibuat gagal bertahun-tahun sekaligus; dan
\emph{pil darurat}, takaran \omterm{def:g8:puberty:hormones}{hormon} susulan dalam hitungan
hari yang bekerja terutama dengan menunda \omterm{prop:g8:reproductive-systems:cycle}{ovulasi} yang belum
terjadi --- rem upaya terakhir, tersedia di apotek, dan
bukan pengganti sebuah cara yang dipakai \emph{sebelumnya}.
Yang tidak dilakukan satu pun darinya: apa pun terhadap
infeksi.
\end{example}

\begin{remark}[Apa yang tidak bekerja]\label{rem:g8:contraception:notwork}
Cara halaman sekolahnya, diukur terhadap rantainya:
``menghitung hari aman'' karam pada keberubahan
\cref{rem:g8:reproductive-systems:living} --- tanggal
\omterm{prop:g8:reproductive-systems:cycle}{ovulasinya} tak dapat diketahui sebelumnya, dan spermanya
bertahan berhari-hari menantikannya; ``menarik diri tepat
waktu'' karam pada kenyataan bahwa spermanya sudah berjalan
sebelum babak akhirnya. Keduanya gagal persis di tempat yang
diramalkan teori mata-rantai-bernamanya: keduanya tidak
menghadang satu mata rantai pun secara andal. Dongeng
bukanlah sebuah cara.
\end{remark}

\section{Pilihan, tanggung jawab, dan peta pertolongannya}

\begin{proposition}[Pilihannya nyata dan dipikul bersama]\label{prop:g8:contraception:choice}
Fakta yang dipaksakan mekanismenya:
\begin{itemize}
  \item \omterm{def:g5:human-reproduction:pregnancy}{kehamilan} membutuhkan satu perbuatan, bukan
        panjangnya sebuah hubungan: rantainya tidak
        menghitung kesempatan;
  \item ia adalah rantai yang \emph{dipikul bersama}: \omterm{def:g5:first-look-at-cells:cell}{sel}
        kedua pasangannya, tanggung jawab kedua pasangannya
        --- \omterm{def:g8:contraception:condom}{kondom} dan pilnya bukanlah urusan satu jenis
        kelamin masing-masing, melainkan sebuah percakapan;
  \item dan tidak ada cara yang sempurna kecuali pantangan,
        walaupun yang baik mendekatinya bila dipakai persis
        seperti rancangannya --- ``dipakai dengan benar''
        itulah tempat sebagian besar kegagalan di dunia
        nyata bersarang.
\end{itemize}
\end{proposition}
\admitted

\begin{example}[Peta pertolongannya]\label{ex:g8:contraception:help}
Alamat pada \cref{ex:g8:puberty:questions}, diperluas:
perawat sekolah dan dokter keluarga menjawab pertanyaan
\omterm{prop:g8:contraception:principle}{kontrasepsi} dengan rahasia terjaga --- termasuk dari anak di
bawah umur; pusat perencanaan keluarga ada di setiap daerah
justru untuk ini, cuma-cuma dan tanpa menghakimi; apoteker
menangani jam pil daruratnya. Sistemnya dibangun supaya
tidak seorang pun harus menyusuri rantainya dengan berbekal
dongeng --- peta pertolongannya adalah bagian dari caranya.
\end{example}

\begin{method}[Menilai cara apa pun yang kamu dengar]\label{met:g8:contraception:evaluate}
Untuk setiap cara yang diklaim, tanyakan berurutan:
\begin{enumerate}
  \item \emph{mata rantai yang mana} pada rantainya yang
        dihadangnya --- dapatkah kamu menyebutnya di petanya?
  \item \emph{seandal apa}, dipakai seperti orang nyata memakainya?
  \item apakah ia melindungi dari \emph{infeksi} juga, atau
        haruskah \omterm{def:g8:contraception:condom}{kondom} menjadi mitranya?
  \item apa yang dituntutnya --- ingatan harian, sebuah
        resep, sebuah pemasangan --- dan apakah itu cocok
        bagi pemakainya?
\end{enumerate}
Tanpa mata rantai yang dapat disebutkan, tidak ada cara
(\cref{rem:g8:contraception:notwork} gagal pada pertanyaan
1).
\end{method}

\section{Latihan}

\begin{exercise}[$\star$]\label{exo:g8:contraception:1}
Nyatakan rantai yang dibutuhkan \omterm{def:g5:human-reproduction:pregnancy}{kehamilan}, lalu rumuskan
\omterm{prop:g8:contraception:principle}{kontrasepsi} terhadapnya.
\end{exercise}

\begin{exercise}[$\star$]\label{exo:g8:contraception:2}
Mata rantai mana yang dihadang \omterm{def:g8:contraception:condom}{kondomnya}? Dan pilnya?
\end{exercise}

\begin{exercise}[$\star$]\label{exo:g8:contraception:3}
Jasa kedua apa yang hanya diberikan \omterm{def:g8:contraception:condom}{kondomnya}?
\end{exercise}

\begin{exercise}[$\star$]\label{exo:g8:contraception:4}
Bagaimana pil daruratnya terutama bekerja, dan ia bukan
apa?
\end{exercise}

\begin{exercise}[$\star$]\label{exo:g8:contraception:5}
Sebutkan tiga alamat pada peta pertolongannya, beserta apa
yang ditawarkan masing-masing.
\end{exercise}

\begin{exercise}[$\star$]\label{exo:g8:contraception:6}
Mengapa ``menghitung hari aman'' tidak andal? Dua alasan
dari \cref{rem:g8:contraception:notwork}.
\end{exercise}

\begin{exercise}[$\star\star$]\label{exo:g8:contraception:7}
Pilahlah rak \cref{ex:g8:contraception:shelf} menurut mata
rantai yang diserangnya, tambahkan \omterm{def:g8:contraception:condom}{kondom} dan pilnya.
\end{exercise}

\begin{exercise}[$\star\star$]\label{exo:g8:contraception:8}
Mengapa dokter begitu sering berkata ``pil \emph{dan}
\omterm{def:g8:contraception:condom}{kondom}'' kepada pasangan muda? Dua perlindungan, satu
kalimat masing-masing.
\end{exercise}

\begin{exercise}[$\star\star$]\label{exo:g8:contraception:9}
Jelaskan, dengan
\cref{prop:g8:reproductive-systems:cycle}, dalam keadaan apa
pilnya menahan daurnya, dan mengapa hal itu mencegah
\omterm{def:g5:human-reproduction:pregnancy}{kehamilan}.
\end{exercise}

\begin{exercise}[$\star\star$]\label{exo:g8:contraception:10}
``\omterm{def:g5:first-look-at-cells:cell}{Sel} kedua pasangannya, urusan kedua pasangannya.''
Pertahankan kalimat itu terhadap kebiasaan menyerahkan
\omterm{prop:g8:contraception:principle}{kontrasepsi} kepada satu sisi pasangannya saja.
\end{exercise}

\begin{exercise}[$\star\star$]\label{exo:g8:contraception:11}
Jalankan \cref{met:g8:contraception:evaluate} pada: (a)
susuk hormonalnya; (b) ``menarik diri tepat waktu''.
\end{exercise}

\begin{exercise}[$\star\star\star$]\label{exo:g8:contraception:12}
``\omterm{prop:g8:contraception:principle}{Kontrasepsi} adalah faal terapan: setiap cara di raknya
adalah sebuah bab buku ini, yang dibidikkan.'' Berikan
alasannya dengan tiga cara dan tiga babnya.
\end{exercise}

\section{Soal: Konsultasi Perencanaan Keluarga}

\begin{problem}[{Soal akhir pekan --- sepetang penuh
pertanyaan di pusat perencanaannya}]\label{pb:g8:contraception:1}
Sepetang (rekaan) seorang penyuluh perencanaan keluarga:
empat pengunjung, empat keadaan, satu peta rantai di
dindingnya.

\medskip\noindent\textbf{Bagian I --- Peta di
dindingnya.}
\begin{enumerate}
  \item Gambarlah peta dindingnya: rantainya dari \omterm{def:g8:sexual-reproduction:gametes}{gamet}
        sampai penanaman, dengan ketiga titik serangan pada
        \cref{prop:g8:contraception:principle} ditandai.
  \item Pengunjung 1 bertanya ``bagaimana satu pil sehari
        dapat menghentikan apa pun?'' Jawablah dengan
        permesinannya: bahasa siapa yang dituturkan pilnya,
        dan kepada siapa?
  \item Pengunjung 1 bertanya lanjut: ``dan bila saya
        melupakan beberapa?'' Tempatkan bahayanya pada
        petanya --- apa yang berjalan lagi ketika isyarat
        \omterm{def:g8:puberty:hormones}{hormonnya} lowong?
  \item Mengapa anjuran \omterm{def:g8:contraception:condom}{kondom} sang penyuluh bertahan
        melewati setiap pilihan cara yang lain? (Tugas
        gandanya.)
\end{enumerate}

\medskip\noindent\textbf{Bagian II --- Keadaannya.}
\begin{enumerate}[resume]
  \item Pengunjung 2, tujuh belas tahun dan cemas, percaya
        bahwa sebuah aplikasi hari aman membuat cara lain
        tidak perlu. Betulkan kepercayaan itu dengan
        lembut: dua fakta biologis mana yang mengalahkan
        penanggalannya?
  \item Pengunjung 3 membutuhkan sesuatu ``tanpa ingatan
        harian''. Tawarkan dua calon dari raknya beserta
        mata rantainya.
  \item Pengunjung 4 datang dalam sehari sesudah sebuah
        \omterm{def:g8:contraception:condom}{kondom} pecah. Sebutkan pilihannya, jamnya,
        mekanisme utamanya --- dan cara tetap yang akan
        dibahas sang penyuluh untuk sesudahnya.
  \item Keempatnya pulang dengan kalimat penutup yang sama
        tentang infeksi. Tuliskanlah.
\end{enumerate}

\medskip\noindent\textbf{Bagian III --- Pembekalan akhirnya.}
\begin{enumerate}[resume]
  \item Sang peserta pelatihan bertanya mengapa pusatnya
        mengajarkan rantainya alih-alih sekadar membagikan
        caranya. Jawablah dengan pertanyaan pertama pada
        \cref{met:g8:contraception:evaluate}.
  \item Pilahlah cara sepetang itu menurut apa yang
        dituntut masing-masing dari pemakainya (pertanyaan 4
        pada caranya).
  \item Sang peserta pelatihan bertanya-tanya mengapa
        kunjungan anak di bawah umur dirahasiakan secara
        sengaja. Jawablah dengan logika
        \cref{ex:g8:contraception:help}.
  \item Tutuplah petangnya dalam satu kalimat: apa yang
        sebenarnya dibagikan pusatnya, di luar kotaknya.
\end{enumerate}
\end{problem}
